% Options for packages loaded elsewhere
% Options for packages loaded elsewhere
\PassOptionsToPackage{unicode}{hyperref}
\PassOptionsToPackage{hyphens}{url}
\PassOptionsToPackage{dvipsnames,svgnames,x11names}{xcolor}
%
\documentclass[
  brazilian,
  letterpaper,
  DIV=11,
  numbers=noendperiod]{scrartcl}
\usepackage{xcolor}
\usepackage{amsmath,amssymb}
\setcounter{secnumdepth}{5}
\usepackage{iftex}
\ifPDFTeX
  \usepackage[T1]{fontenc}
  \usepackage[utf8]{inputenc}
  \usepackage{textcomp} % provide euro and other symbols
\else % if luatex or xetex
  \usepackage{unicode-math} % this also loads fontspec
  \defaultfontfeatures{Scale=MatchLowercase}
  \defaultfontfeatures[\rmfamily]{Ligatures=TeX,Scale=1}
\fi
\usepackage{lmodern}
\ifPDFTeX\else
  % xetex/luatex font selection
\fi
% Use upquote if available, for straight quotes in verbatim environments
\IfFileExists{upquote.sty}{\usepackage{upquote}}{}
\IfFileExists{microtype.sty}{% use microtype if available
  \usepackage[]{microtype}
  \UseMicrotypeSet[protrusion]{basicmath} % disable protrusion for tt fonts
}{}
\makeatletter
\@ifundefined{KOMAClassName}{% if non-KOMA class
  \IfFileExists{parskip.sty}{%
    \usepackage{parskip}
  }{% else
    \setlength{\parindent}{0pt}
    \setlength{\parskip}{6pt plus 2pt minus 1pt}}
}{% if KOMA class
  \KOMAoptions{parskip=half}}
\makeatother
% Make \paragraph and \subparagraph free-standing
\makeatletter
\ifx\paragraph\undefined\else
  \let\oldparagraph\paragraph
  \renewcommand{\paragraph}{
    \@ifstar
      \xxxParagraphStar
      \xxxParagraphNoStar
  }
  \newcommand{\xxxParagraphStar}[1]{\oldparagraph*{#1}\mbox{}}
  \newcommand{\xxxParagraphNoStar}[1]{\oldparagraph{#1}\mbox{}}
\fi
\ifx\subparagraph\undefined\else
  \let\oldsubparagraph\subparagraph
  \renewcommand{\subparagraph}{
    \@ifstar
      \xxxSubParagraphStar
      \xxxSubParagraphNoStar
  }
  \newcommand{\xxxSubParagraphStar}[1]{\oldsubparagraph*{#1}\mbox{}}
  \newcommand{\xxxSubParagraphNoStar}[1]{\oldsubparagraph{#1}\mbox{}}
\fi
\makeatother


\usepackage{longtable,booktabs,array}
\newcounter{none} % for unnumbered tables
\usepackage{calc} % for calculating minipage widths
% Correct order of tables after \paragraph or \subparagraph
\usepackage{etoolbox}
\makeatletter
\patchcmd\longtable{\par}{\if@noskipsec\mbox{}\fi\par}{}{}
\makeatother
% Allow footnotes in longtable head/foot
\IfFileExists{footnotehyper.sty}{\usepackage{footnotehyper}}{\usepackage{footnote}}
\makesavenoteenv{longtable}
\usepackage{graphicx}
\makeatletter
\newsavebox\pandoc@box
\newcommand*\pandocbounded[1]{% scales image to fit in text height/width
  \sbox\pandoc@box{#1}%
  \Gscale@div\@tempa{\textheight}{\dimexpr\ht\pandoc@box+\dp\pandoc@box\relax}%
  \Gscale@div\@tempb{\linewidth}{\wd\pandoc@box}%
  \ifdim\@tempb\p@<\@tempa\p@\let\@tempa\@tempb\fi% select the smaller of both
  \ifdim\@tempa\p@<\p@\scalebox{\@tempa}{\usebox\pandoc@box}%
  \else\usebox{\pandoc@box}%
  \fi%
}
% Set default figure placement to htbp
\def\fps@figure{htbp}
\makeatother


% definitions for citeproc citations
\NewDocumentCommand\citeproctext{}{}
\NewDocumentCommand\citeproc{mm}{%
  \begingroup\def\citeproctext{#2}\cite{#1}\endgroup}
\makeatletter
 % allow citations to break across lines
 \let\@cite@ofmt\@firstofone
 % avoid brackets around text for \cite:
 \def\@biblabel#1{}
 \def\@cite#1#2{{#1\if@tempswa , #2\fi}}
\makeatother
\newlength{\cslhangindent}
\setlength{\cslhangindent}{1.5em}
\newlength{\csllabelwidth}
\setlength{\csllabelwidth}{3em}
\newenvironment{CSLReferences}[2] % #1 hanging-indent, #2 entry-spacing
 {\begin{list}{}{%
  \setlength{\itemindent}{0pt}
  \setlength{\leftmargin}{0pt}
  \setlength{\parsep}{0pt}
  % turn on hanging indent if param 1 is 1
  \ifodd #1
   \setlength{\leftmargin}{\cslhangindent}
   \setlength{\itemindent}{-1\cslhangindent}
  \fi
  % set entry spacing
  \setlength{\itemsep}{#2\baselineskip}}}
 {\end{list}}
\usepackage{calc}
\newcommand{\CSLBlock}[1]{\hfill\break\parbox[t]{\linewidth}{\strut\ignorespaces#1\strut}}
\newcommand{\CSLLeftMargin}[1]{\parbox[t]{\csllabelwidth}{\strut#1\strut}}
\newcommand{\CSLRightInline}[1]{\parbox[t]{\linewidth - \csllabelwidth}{\strut#1\strut}}
\newcommand{\CSLIndent}[1]{\hspace{\cslhangindent}#1}

\ifLuaTeX
\usepackage[bidi=basic,shorthands=off]{babel}
\else
\usepackage[bidi=default,shorthands=off]{babel}
\fi
\ifLuaTeX
  \usepackage{selnolig} % disable illegal ligatures
\fi


\setlength{\emergencystretch}{3em} % prevent overfull lines

\providecommand{\tightlist}{%
  \setlength{\itemsep}{0pt}\setlength{\parskip}{0pt}}



 


\KOMAoption{captions}{tableheading}
\makeatletter
\@ifpackageloaded{tcolorbox}{}{\usepackage[skins,breakable]{tcolorbox}}
\@ifpackageloaded{fontawesome5}{}{\usepackage{fontawesome5}}
\definecolor{quarto-callout-color}{HTML}{909090}
\definecolor{quarto-callout-note-color}{HTML}{0758E5}
\definecolor{quarto-callout-important-color}{HTML}{CC1914}
\definecolor{quarto-callout-warning-color}{HTML}{EB9113}
\definecolor{quarto-callout-tip-color}{HTML}{00A047}
\definecolor{quarto-callout-caution-color}{HTML}{FC5300}
\definecolor{quarto-callout-color-frame}{HTML}{acacac}
\definecolor{quarto-callout-note-color-frame}{HTML}{4582ec}
\definecolor{quarto-callout-important-color-frame}{HTML}{d9534f}
\definecolor{quarto-callout-warning-color-frame}{HTML}{f0ad4e}
\definecolor{quarto-callout-tip-color-frame}{HTML}{02b875}
\definecolor{quarto-callout-caution-color-frame}{HTML}{fd7e14}
\makeatother
\makeatletter
\@ifpackageloaded{caption}{}{\usepackage{caption}}
\AtBeginDocument{%
\ifdefined\contentsname
  \renewcommand*\contentsname{Índice}
\else
  \newcommand\contentsname{Índice}
\fi
\ifdefined\listfigurename
  \renewcommand*\listfigurename{Lista de Figuras}
\else
  \newcommand\listfigurename{Lista de Figuras}
\fi
\ifdefined\listtablename
  \renewcommand*\listtablename{Lista de Tabelas}
\else
  \newcommand\listtablename{Lista de Tabelas}
\fi
\ifdefined\figurename
  \renewcommand*\figurename{Figura}
\else
  \newcommand\figurename{Figura}
\fi
\ifdefined\tablename
  \renewcommand*\tablename{Tabela}
\else
  \newcommand\tablename{Tabela}
\fi
}
\@ifpackageloaded{float}{}{\usepackage{float}}
\floatstyle{ruled}
\@ifundefined{c@chapter}{\newfloat{codelisting}{h}{lop}}{\newfloat{codelisting}{h}{lop}[chapter]}
\floatname{codelisting}{Listagem}
\newcommand*\listoflistings{\listof{codelisting}{Lista de Listagens}}
\makeatother
\makeatletter
\makeatother
\makeatletter
\@ifpackageloaded{caption}{}{\usepackage{caption}}
\@ifpackageloaded{subcaption}{}{\usepackage{subcaption}}
\makeatother
\usepackage{bookmark}
\IfFileExists{xurl.sty}{\usepackage{xurl}}{} % add URL line breaks if available
\urlstyle{same}
\hypersetup{
  pdftitle={Exercício 1: Biodiversidade{,} níveis de organização e montagem de comunidades aplicados à restauração ecológica},
  pdflang={pt-BR},
  colorlinks=true,
  linkcolor={blue},
  filecolor={Maroon},
  citecolor={Blue},
  urlcolor={Blue},
  pdfcreator={LaTeX via pandoc}}


\title{Exercício 1: Biodiversidade, níveis de organização e montagem de
comunidades aplicados à restauração ecológica}
\usepackage{etoolbox}
\makeatletter
\providecommand{\subtitle}[1]{% add subtitle to \maketitle
  \apptocmd{\@title}{\par {\large #1 \par}}{}{}
}
\makeatother
\subtitle{Ecologia da Restauração --- UFPE, Departamento de Botânica}
\author{}
\date{}
\begin{document}
\maketitle


\section{Cenário}\label{cenuxe1rio}

\pandocbounded{\includegraphics[keepaspectratio]{exercicio_restauracao_vellend_files/mediabag/image_processing2021.jpg}}

\pandocbounded{\includegraphics[keepaspectratio]{exercicio_restauracao_vellend_files/mediabag/ft-1.jpg}}

Um assentamento de reforma agrária foi implantado numa antiga usina de
cana-de-açúcar no agreste/zona da mata de Pernambuco. A paisagem
resultante é um mosaico de fragmentos florestais de Mata Atlântica em
diferentes estágios de regeneração, cercados por talhões abandonados de
cana, pastagens degradadas e pequenas roças de subsistência. A
associação de assentados, em parceria com uma ONG ambiental e a
universidade, decidiu iniciar um programa de restauração ecológica com
dois objetivos declarados:

\begin{enumerate}
\def\labelenumi{\arabic{enumi}.}
\tightlist
\item
  \textbf{Reconectar} os fragmentos florestais remanescentes por meio de
  corredores de vegetação nativa;
\item
  \textbf{Proteger} as nascentes e as matas ciliares dos riachos que
  cortam o assentamento, hoje amplamente desmatadas e assoreadas.
\end{enumerate}

O plano de restauração prevê ações distintas conforme o contexto:
plantio de mudas em áreas totalmente descobertas, condução da
regeneração natural em áreas com banco de sementes e rebrota
preservados, e enriquecimento de fragmentos já em recuperação espontânea
há cerca de 15 anos.

Você foi contratado(a) como consultor(a) técnico(a) para ajudar a equipe
a interpretar o processo ecológico em curso e a justificar
cientificamente as decisões de manejo.

\begin{tcolorbox}[enhanced jigsaw, arc=.35mm, bottomrule=.15mm, bottomtitle=1mm, breakable, colback=white, colbacktitle=quarto-callout-note-color!10!white, colframe=quarto-callout-note-color-frame, coltitle=black, left=2mm, leftrule=.75mm, opacityback=0, opacitybacktitle=0.6, rightrule=.15mm, title=\textcolor{quarto-callout-note-color}{\faInfo}\hspace{0.5em}{Formato do exercício}, titlerule=0mm, toprule=.15mm, toptitle=1mm]

Exercício domiciliar, individual, com consulta livre a bibliografia,
anotações de aula e à internet. O que se avalia não é a capacidade de
encontrar a ``resposta certa'', mas a capacidade de \textbf{aplicar
corretamente os conceitos ao cenário específico}, com raciocínio
explícito e articulado às particularidades da Mata Atlântica nordestina
e do contexto socioambiental do assentamento.

\end{tcolorbox}

\section{Objetivos de aprendizagem}\label{objetivos-de-aprendizagem}

Ao final deste exercício, espera-se que o(a) estudante seja capaz de:

\begin{itemize}
\tightlist
\item
  Diferenciar e aplicar corretamente os níveis de organização biológica
  (população, comunidade, paisagem) a uma situação real de manejo;
\item
  Caracterizar biodiversidade em múltiplas dimensões (taxonômica,
  funcional, filogenética) e discutir por que essas dimensões podem
  responder de forma diferente à restauração;
\item
  Aplicar o arcabouço teórico de Vellend (2010) --- \textbf{seleção,
  deriva ecológica, dispersão e especiação/mutação} --- para explicar a
  montagem de comunidades ao longo da sucessão;
\item
  Relacionar os quatro processos de Vellend às decisões práticas de
  restauração descritas no cenário (plantio ativo vs.~regeneração
  natural vs.~enriquecimento, corredores, proteção de nascentes).
\end{itemize}

\section{Parte 1 --- Biodiversidade e níveis de organização
biológica}\label{parte-1-biodiversidade-e-nuxedveis-de-organizauxe7uxe3o-bioluxf3gica}

\textbf{1.1.} No fragmento mais antigo (regeneração espontânea de
\textasciitilde15 anos), um levantamento florístico registrou 40
espécies lenhosas. Num trecho de mata ciliar ainda degradada, apenas 8
espécies foram registradas, mas duas delas são endêmicas de matas
ciliares nordestinas e não ocorrem em nenhum outro ponto do
assentamento.

Discuta: usar apenas a \textbf{riqueza de espécies} para comparar esses
dois pontos e decidir ``onde investir recursos de restauração'' seria
uma decisão ecologicamente bem informada? Incorpore em sua resposta pelo
menos duas dimensões adicionais de biodiversidade (por exemplo,
diversidade funcional, filogenética, beta-diversidade ou diversidade
genética) e explique o que cada uma agregaria à decisão.

\textbf{1.2.} Explique, usando exemplos concretos do cenário, a
diferença entre analisar o processo de restauração no nível de
\textbf{população} (ex.: uma população de uma espécie pioneira
específica) e no nível de \textbf{comunidade} (ex.: o conjunto de
espécies lenhosas de um fragmento). Por que um plano de restauração de
paisagem não pode se apoiar em apenas um desses níveis?

\section{Parte 2 --- Processos de montagem de comunidades (Vellend
2010)}\label{parte-2-processos-de-montagem-de-comunidades-vellend-2010}

Vellend (2010) propõe que toda a dinâmica de comunidades ecológicas pode
ser explicada por quatro processos de alto nível: \textbf{seleção}
(determinística, ligada a diferenças de aptidão entre espécies em
resposta ao ambiente biótico e abiótico, incluindo interações),
\textbf{deriva ecológica} (mudanças estocásticas em abundância não
ligadas a diferenças de aptidão), \textbf{dispersão} (movimento de
organismos/propágulos no espaço) e \textbf{especiação} (origem de nova
variação genética/de espécies).

\textbf{2.1.} Para cada um dos quatro processos, identifique \textbf{um
elemento concreto do cenário} do assentamento em que esse processo
provavelmente está atuando (ou deveria ser manejado) e justifique
brevemente. Não vale apenas citar a definição --- é preciso ancorar no
cenário.

\textbf{2.2.} Compare os dois pontos de partida propostos no plano de
restauração:

\begin{itemize}
\tightlist
\item
  \begin{enumerate}
  \def\labelenumi{(\alph{enumi})}
  \tightlist
  \item
    uma área de solo exposto, ex-canavial, sem banco de sementes viável,
    onde será feito plantio de mudas de várias espécies nativas;
  \end{enumerate}
\item
  \begin{enumerate}
  \def\labelenumi{(\alph{enumi})}
  \setcounter{enumi}{1}
  \tightlist
  \item
    um trecho de mata ciliar degradada, mas com banco de sementes e
    capacidade de rebrota de tocos ainda presentes, onde a equipe optou
    por conduzir a regeneração natural.
  \end{enumerate}
\end{itemize}

Discuta como o \textbf{peso relativo} dos quatro processos de Vellend
difere entre (a) e (b) na fase inicial da restauração. Em qual dos dois
cenários a \textbf{dispersão} tende a ser mais limitante, e por quê? Que
implicação prática isso tem para o desenho da intervenção (por exemplo,
necessidade de plantio, poleiros artificiais, translocação de
serapilheira, etc.)?

\section{Parte 3 --- Sucessão ecológica e corredores na
paisagem}\label{parte-3-sucessuxe3o-ecoluxf3gica-e-corredores-na-paisagem}

\textbf{3.1.} Ao longo da sucessão secundária de uma área que era
canavial (portanto, com histórico de perturbação intensa, uso de
agrotóxicos e possível esgotamento/compactação do solo), descreva como
você espera que o \textbf{peso relativo dos quatro processos de Vellend
mude} entre o ano 1, o ano 15 e o ano 50 após o início da restauração.
Sua resposta deve deixar claro \emph{por que} esse peso muda (por
exemplo: o que muda no ambiente, na comunidade já estabelecida, ou na
paisagem ao redor, que altera a importância relativa de seleção, deriva,
dispersão e especiação).

\textbf{3.2.} A equipe técnica propôs conectar fragmentos via corredores
ripários ao longo dos riachos, em vez de plantar blocos florestais
isolados no meio das antigas áreas de cana. Do ponto de vista da
\textbf{dispersão} como processo de montagem de comunidades, que
vantagem ecológica essa escolha de desenho espacial oferece para a
paisagem como um todo, em comparação com fragmentos isolados sem
conectividade?

\textbf{3.3.} Um dos assentados questiona por que a equipe insiste tanto
em proteger as nascentes e não apenas plantar árvores ``onde há mais
espaço disponível''. Elabore uma resposta que combine (i) um argumento
ecológico ligado aos processos de montagem de comunidades ou aos níveis
de organização biológica discutidos neste exercício, e (ii) um argumento
ligado à função hidrológica/ambiental das nascentes e matas ciliares na
paisagem. A resposta deve ser compreensível para um público não
especialista, sem perder rigor conceitual.

\section{Formato de entrega}\label{formato-de-entrega}

\begin{itemize}
\tightlist
\item
  Documento de texto (PDF), respostas identificadas por item (1.1, 1.2,
  2.1 etc.);
\item
  Extensão sugerida: 150--300 palavras por item --- respostas mais
  longas não são necessariamente melhores; respostas cirúrgicas e bem
  ancoradas no cenário valem mais que respostas genéricas de manual;
\item
  Cite ao menos Vellend (2010) e mais uma referência complementar de sua
  escolha, relacionada a sucessão ecológica ou restauração na Mata
  Atlântica/Nordeste;
\item
  É permitida consulta a qualquer fonte, incluindo ferramentas de IA ---
  mas respostas genéricas, que não usem os elementos específicos do
  cenário (canavial, nascentes, assentamento, corredores ripários),
  perderão pontos, pois o exercício foi desenhado exatamente para exigir
  essa aplicação situada.
\end{itemize}

\begin{tcolorbox}[enhanced jigsaw, arc=.35mm, bottomrule=.15mm, bottomtitle=1mm, breakable, colback=white, colbacktitle=quarto-callout-important-color!10!white, colframe=quarto-callout-important-color-frame, coltitle=black, left=2mm, leftrule=.75mm, opacityback=0, opacitybacktitle=0.6, rightrule=.15mm, title=\textcolor{quarto-callout-important-color}{\faExclamation}\hspace{0.5em}{Guia de correção --- o que se espera nas respostas (uso do
professor/monitor)}, titlerule=0mm, toprule=.15mm, toptitle=1mm]

\textbf{1.1} --- Espera-se que o(a) estudante rejeite a riqueza de
espécies como critério único, reconhecendo que ela é cega a raridade,
endemismo e função. Resposta forte menciona explicitamente que
diversidade funcional captura papéis ecológicos (ex.: fixação de N,
dispersão de sementes, resistência a seca) que não são visíveis apenas
contando espécies; diversidade filogenética/genética capta singularidade
evolutiva (relevante para as duas endêmicas); beta-diversidade entre os
dois pontos indicaria se eles são complementares na paisagem. Deve
concluir que os 8-espécies-endêmicas podem ter prioridade de conservação
mesmo com riqueza baixa.

\textbf{1.2} --- Espera-se articulação clara: nível de população =
dinâmica de uma espécie específica (recrutamento, mortalidade, estrutura
etária de uma pioneira, por exemplo); nível de comunidade = composição e
interações entre múltiplas espécies coexistindo num fragmento. Resposta
forte explica que decisões de manejo no nível populacional (ex.: quantas
mudas de uma espécie plantar) não garantem, por si só, uma comunidade
funcional e resiliente --- e vice-versa, dados de comunidade agregados
podem esconder o colapso de uma população-chave (ex.: uma espécie
hidrófila essencial para a mata ciliar).

\textbf{2.1} --- Respostas aceitáveis variam, mas devem ancorar no
cenário: \emph{seleção} --- filtros ambientais do solo
compactado/exposto do ex-canavial favorecendo pioneiras tolerantes a sol
e estresse hídrico; \emph{deriva} --- em manchas pequenas e
recém-plantadas, flutuações estocásticas de mortalidade podem eliminar
espécies raras independentemente de sua aptidão; \emph{dispersão} ---
fragmentos isolados sem corredores dependem de dispersores (aves,
morcegos) que atravessam a matriz de pasto/cana, algo limitado;
\emph{especiação/mutação} --- processo de longo prazo, pouco relevante
na escala de décadas do projeto, mas pode ser mencionado via
variabilidade genética das mudas usadas no plantio (evitar gargalos
genéticos).

\textbf{2.2} --- Espera-se reconhecimento de que em (a) --- solo
exposto, sem banco de sementes --- a dispersão e a seleção assumem peso
desproporcional desde o início: praticamente toda a composição inicial
depende de propágulos trazidos de fora (plantio ativo é, na prática,
dispersão assistida) e do que sobrevive aos filtros abióticos severos.
Em (b), com banco de sementes e rebrota, a comunidade ``herda'' parte da
composição pré-existente, e a deriva/seleção atuam sobre esse pool já
presente, reduzindo a dependência crítica de dispersão externa. Resposta
forte conclui que (a) é onde a dispersão é mais limitante e justifica
intervenções ativas (plantio, poleiros, corredores) com mais urgência.

\textbf{3.1} --- Resposta forte descreve uma trajetória: ano 1 ---
seleção abiótica dominante (estresse do solo ex-canavial, alta
luminosidade) e deriva forte por baixa densidade populacional; dispersão
crítica e limitante por isolamento de fragmentos-fonte; ano 15 --- com
dossel formado, seleção biótica (competição, facilitação) ganha peso
sobre seleção abiótica; deriva perde importância relativa à medida que
populações crescem; dispersão ainda relevante para incorporar espécies
tardias/zoocóricas; ano 50 --- comunidade mais próxima de equilíbrio
dinâmico, seleção biótica (interações, incluindo herbivoria e patógenos
denso-dependentes) mais proeminente, especiação ainda irrelevante na
escala temporal. Não há uma única resposta ``correta'', mas a lógica
causal deve ser explícita.

\textbf{3.2} --- Espera-se o argumento de que corredores ripários
funcionam como estruturas físicas que aumentam a conectividade funcional
da paisagem, reduzindo a limitação por dispersão entre fragmentos ---
favorecendo fluxo gênico, recolonização após perturbações locais e
movimento de dispersores de sementes (fauna). Contraste explícito com
fragmentos isolados, onde a dispersão eficaz depende de a fauna
atravessar matriz hostil (pasto, cana), algo raro para muitas espécies
de sub-bosque com baixa vagilidade.

\textbf{3.3} --- Espera-se resposta híbrida: argumento ecológico
(nascentes/matas ciliares protegem cabeceiras de drenagem que sustentam
banco de sementes, umidade e conectividade estrutural --- ligando de
volta a Vellend, são pontos-chave de dispersão e refúgios de baixa
perturbação que ancoram a seleção biótica ao redor) + argumento
hidrológico (proteção de vazão de nascentes, controle de assoreamento,
manutenção de qualidade de água para o próprio assentamento). Respostas
que ignorem um dos dois eixos são incompletas.

\end{tcolorbox}

\section*{Referências}\label{referuxeancias}
\addcontentsline{toc}{section}{Referências}

\protect\phantomsection\label{refs}
\begin{CSLReferences}{1}{1}
\bibitem[\citeproctext]{ref-vellend2010}
Vellend, Mark. 2010. {``Conceptual Synthesis in Community Ecology''}.
\emph{The Quarterly Review of Biology} 85 (2): 183--206.
\url{https://doi.org/10.1086/652373}.

\end{CSLReferences}




\end{document}
